% !TEX program = xelatex
\documentclass[11pt,a4paper]{article}
\usepackage{algo-preamble}

\title{\textbf{L01 \textperiodcentered{} Εισαγωγικά}}
\author{Algorithms Class Hub}
\date{\small Αλγόριθμοι και Πολυπλοκότητα (K17) \textperiodcentered{} ΕΚΠΑ}

\begin{document}
\maketitle

\begin{center}\itshape\small
Τι είναι αλγόριθμος, πώς μετράμε την «καλή» επίδοσή του, και ποια θα είναι η ροή
του εξαμήνου. Στιγμιότυπο, διάσταση, στοιχειώδης πράξη, τρεις τύποι πολυπλοκότητας.
\end{center}

\section{Καλώς ήρθες}

Αυτή είναι η αφετηρία για όλο το μάθημα. Στόχος της διάλεξης: μέχρι το τέλος, να
ξέρεις \textbf{τι ακριβώς εννοούμε όταν λέμε «αλγόριθμος»}, με ποιο εργαλείο τον
συγκρίνουμε με άλλους, και ποια διαδρομή θα ακολουθήσουμε από εδώ μέχρι την
εξεταστική.

\begin{intuition}
\textbf{Πρακτικά:} ένας αλγόριθμος είναι μια συνταγή. Μια συνταγή μαγειρικής σου
λέει «βάλε 2 αυγά, ανακάτεψε για 30 δευτερόλεπτα» --- βήματα ξεκάθαρα, σε σειρά,
και μετά από πεπερασμένη ώρα έχεις φαγητό. Ένας αλγόριθμος κάνει το ίδιο για ένα
\textbf{υπολογιστικό πρόβλημα}.
\end{intuition}

\section{Από πού έρχεται η λέξη «αλγόριθμος»}

Από τον \textbf{Al-Khwarizmi} --- Πέρση μαθηματικό του 9ου αιώνα. Έγραψε ένα βιβλίο
για μεθοδικές διαδικασίες λύσης εξισώσεων· όταν το έργο του μεταφράστηκε στα
Λατινικά, το όνομά του παραφθάρηκε σε \textit{Algoritmi}, και η λέξη επιβίωσε.

\section{Γιατί χάνουμε χρόνο σε ένα μάθημα Αλγορίθμων;}

Εύλογο ερώτημα --- ειδικά τώρα που υπάρχουν LLMs που λύνουν ασκήσεις. Δύο
απαντήσεις:

\begin{enumerate}
  \item \textbf{Τα παραδείγματα είναι παντού.} Κοινωνικά δίκτυα (ποιους φίλους να
  σου προτείνει το Instagram), επικοινωνία drones (σε ποιο συντομότερο μονοπάτι να
  πάει το επόμενο), εύρεση διαδρομών (Google Maps), συναλλαγές (κρυπτογραφία).
  Κάτω από όλα αυτά \textbf{κάποιος αλγόριθμος} τρέχει, και η ποιότητά του
  καθορίζει αν η εφαρμογή δουλεύει σε εκατομμύρια χρήστες ή κολλάει στα 100.
  \item \textbf{Τα LLMs σου λένε \textit{τι}, όχι \textit{γιατί}.} Όταν στην
  εξέταση δεις μια άσκηση τύπου «σχεδίασε αλγόριθμο που\dots» θα πρέπει να
  αναγνωρίσεις \textit{ποιο σχήμα} εφαρμόζεται, να αποδείξεις ότι δουλεύει, και να
  αναλύσεις την πολυπλοκότητα. Αυτό το μάθημα είναι ακριβώς αυτή η εξάσκηση.
\end{enumerate}

\section{Οι τρεις στόχοι του εξαμήνου}

Όλα όσα κάνουμε καταλήγουν σε ένα από αυτά τα τρία:

\begin{enumerate}
  \item \textbf{Αποτίμηση} --- δοθέντος ενός αλγορίθμου, πόσο χρόνο και μνήμη
  χρειάζεται;
  \item \textbf{Σύγκριση} --- δοθέντων δύο αλγορίθμων που λύνουν το ίδιο πρόβλημα,
  ποιος είναι καλύτερος και υπό ποιες προϋποθέσεις;
  \item \textbf{Σχεδίαση} --- δοθέντος ενός προβλήματος που δεν έχουμε ξαναδεί,
  \textbf{πώς} στήνουμε έναν καλό αλγόριθμο για αυτό;
\end{enumerate}

\section{Δύο σημαντικές ερωτήσεις πριν προχωρήσουμε}

\subsection{Είναι όλα τα προβλήματα επιλύσιμα από έναν αλγόριθμο;}

\textbf{Όχι.} Υπάρχουν προβλήματα που είναι αποδεδειγμένα \textbf{μη επιλύσιμα} από
οποιονδήποτε αλγόριθμο (το διάσημο παράδειγμα: το \textit{halting problem} ---
«δοθέντος ενός προγράμματος, θα τερματίσει;»). Στο μάθημα δεν θα μπούμε σε
υπολογιστική μη-αποφασισιμότητα, αλλά αξίζει να ξέρεις ότι υπάρχει αυτό το όριο.

\subsection{Υπάρχει για όλα τα επιλύσιμα προβλήματα ένας «καλός» αλγόριθμος;}

Αυτή είναι η ερώτηση που γέννησε ολόκληρο τον τομέα της \textbf{πολυπλοκότητας}.
Υπάρχουν προβλήματα \textbf{επιλύσιμα μεν}, αλλά κανείς δεν ξέρει αν υπάρχει
«γρήγορος» ($=$πολυωνυμικός) αλγόριθμος για αυτά. Το αν $P = NP$ είναι ένα από τα
7 «Millennium problems» με 1 εκ. δολάρια έπαθλο.

\begin{keypoint}
\textbf{Πρόγραμμα και Αλγόριθμος.} Κάθε πρόγραμμα υλοποιεί κάποιον αλγόριθμο. Το
αντίστροφο όχι πάντα ευθέως: ένας αλγόριθμος μπορεί να περιγραφεί σε ψευδογλώσσα,
αλλά για να τρέξει χρειάζεται προγραμματισμό σε συγκεκριμένη γλώσσα. \textbf{Για το
μάθημα τους θεωρούμε ισοδύναμους:} Αλγόριθμος $\leftrightarrow$ Πρόγραμμα.
\end{keypoint}

\section{Παράδειγμα: η ψευδογλώσσα μας}

Πριν μιλήσουμε για πολυπλοκότητα, ας δούμε πώς γράφουμε έναν αλγόριθμο
\textbf{σε ψευδογλώσσα}. Εδώ είναι η δυαδική αναζήτηση σε ταξινομημένο πίνακα
(π.χ.\ τηλεφωνικός κατάλογος):

\begin{codebox}
\begin{verbatim}
Αλγόριθμος ΔυαδικήΑναζήτηση(Π, x):
  Επανέλαβε:
    Σύγκρινε το x με το μεσαίο στοιχείο του Π
    Αν είναι μικρότερο: συνέχισε στο πρώτο ήμισυ
    Αλλιώς: συνέχισε στο δεύτερο ήμισυ
  Έως (x == μεσαίο στοιχείο) ή (ο πίνακας είναι κενός)
\end{verbatim}
\end{codebox}

Δες πώς δουλεύει σε ένα μικρό παράδειγμα --- αναζήτηση του $23$ σε ταξινομημένο
πίνακα $8$ στοιχείων, $[5,\ 11,\ 14,\ 17,\ 23,\ 28,\ 33,\ 41]$:

\begin{itemize}
  \item \textbf{Βήμα 1:} μεσαίο στοιχείο $=17$. Επειδή $23 > 17$, η αναζήτηση
  συνεχίζει στο δεξί μισό $[23,\ 28,\ 33,\ 41]$.
  \item \textbf{Βήμα 2:} μεσαίο στοιχείο $=33$. Επειδή $23 < 33$, η αναζήτηση
  συνεχίζει στο αριστερό μισό $[23,\ 28]$.
  \item \textbf{Βήμα 3:} μεσαίο στοιχείο $=23$. Ισότητα --- το στοιχείο βρέθηκε.
\end{itemize}

Σε πίνακα $8$ στοιχείων χρειάστηκαν \textbf{3 συγκρίσεις}. Σε πίνακα
$1{.}000{.}000$ στοιχείων θα χρειαστούν περίπου \textbf{20}. Αυτή η συμπεριφορά
«κάθε βήμα κόβει στη μέση» είναι ο πυρήνας του τι μελετάμε στο επόμενο μάθημα, την
ασυμπτωτική ανάλυση, και επιστρέφει αργότερα στο Divide \& Conquer.

\section{Πότε ένα πρόγραμμα είναι «χρήσιμο»;}

Σύμφωνα με τη διάλεξη:

\begin{quoteblock}
Ένα πρόγραμμα είναι χρήσιμο όταν, \textbf{σε λογικό χρόνο} και \textbf{σε λογικό
χώρο μνήμης}, εξάγει το \textbf{σωστό αποτέλεσμα}.
\end{quoteblock}

Αυτή η μία πρόταση κρύβει τρεις απαιτήσεις:

\begin{center}
\begin{tabular}{@{}p{3.2cm} p{9.5cm}@{}}
\toprule
\textbf{Απαίτηση} & \textbf{Τι μετράμε} \\
\midrule
\textbf{Ορθότητα} & Το αποτέλεσμα είναι αυτό που υποσχέθηκε ο αλγόριθμος; \\
\textbf{Χρόνος}   & Πόσες στοιχειώδεις πράξεις χρειάζονται; \\
\textbf{Μνήμη}    & Πόσο επιπλέον χώρο καταναλώνει; \\
\bottomrule
\end{tabular}
\end{center}

Οι δύο τελευταίες μαζί αποτελούν την \textbf{πολυπλοκότητα} του αλγορίθμου.

\section{Πρόβλημα, στιγμιότυπο, διάσταση --- τρία διαφορετικά πράγματα}

Είναι εύκολο να τα μπερδέψεις, αλλά κάθε άσκηση εξετάσεων θα σου πει ξεκάθαρα και
τα τρία. Πάμε με δύο παραδείγματα.

\begin{examplebox}
\textbf{Πρόβλημα 1 --- Ταξινόμηση}
\begin{itemize}
  \item \textit{Είσοδος:} $n$ ακέραιοι $a_1, a_2, \dots, a_n$.
  \item \textit{Ζητούμενο:} να ταξινομήσουμε τους ακέραιους σε αύξουσα τάξη.
\end{itemize}

\textbf{Στιγμιότυπο} ενός τέτοιου προβλήματος είναι π.χ.\ $\{8, 3, 5, 1, 9\}$ ---
ένας συγκεκριμένος πίνακας $5$ ακεραίων.

\textbf{Διάσταση} του στιγμιοτύπου είναι το $n = 5$ --- δηλαδή ο αριθμός των
στοιχείων.
\end{examplebox}

\begin{examplebox}
\textbf{Πρόβλημα 2 --- Knapsack (σακίδιο)}
\begin{itemize}
  \item \textit{Είσοδος:} $n$ αντικείμενα με αξίες $c[i]$ και όγκους $w[i]$, και
  ένα σακίδιο χωρητικότητας $b$.
  \item \textit{Ζητούμενο:} να επιλέξουμε σύνολο αντικειμένων με τη μεγαλύτερη
  αξία που να χωράνε στο σακίδιο.
\end{itemize}

\textbf{Στιγμιότυπο:} ένα συγκεκριμένο σύνολο αντικειμένων με συγκεκριμένες
αξίες/όγκους και συγκεκριμένη χωρητικότητα.

\textbf{Διάσταση:} το $n$ --- δηλαδή πάλι ο αριθμός των αντικειμένων.
\end{examplebox}

\begin{keypoint}
\textbf{Διάκριση} --- \textit{Πρόβλημα} $=$ γενική περιγραφή. \textit{Στιγμιότυπο}
$=$ συγκεκριμένα δεδομένα ενός παραδείγματος. \textit{Διάσταση} $=$ αριθμητική
παράμετρος που μετράει το «μέγεθος» του στιγμιοτύπου. Πάντα η πολυπλοκότητα
εκφράζεται ως \textbf{συνάρτηση της διάστασης}.
\end{keypoint}

\section{Πώς μετράμε την πολυπλοκότητα}

\begin{itemize}
  \item \textbf{Μονάδα μέτρησης:} η \textbf{στοιχειώδης (ουσιώδης) πράξη} --- μια
  ενέργεια που υποθέτουμε ότι γίνεται σε σταθερό χρόνο. Παραδείγματα: μία
  πρόσθεση, μία σύγκριση, μία εκχώρηση σε μεταβλητή.
  \item \textbf{Συνάρτηση της διάστασης:} μετράμε \textit{πόσες} στοιχειώδεις
  πράξεις γίνονται ως συνάρτηση του $n$.
\end{itemize}

Η συνάρτηση αυτή είναι ο \textbf{χρόνος εκτέλεσης} του αλγορίθμου, και θα μάθουμε
στο επόμενο μάθημα πώς τη συγκρίνουμε με άλλες χρησιμοποιώντας τα
$O, \Theta, \Omega$.

\section{Τρεις τύποι πολυπλοκότητας}

Όταν λέμε «πόσο χρόνο κάνει ο αλγόριθμος;», για ποιο στιγμιότυπο μιλάμε;
Διαφορετικά στιγμιότυπα της ίδιας διάστασης μπορεί να απαιτούν διαφορετικό χρόνο.
Διακρίνουμε τρεις τύπους:

\begin{itemize}
  \item \textbf{Βέλτιστη περίπτωση} ($C_{\beta\pi}$) --- το «πιο εύκολο»
  στιγμιότυπο.
  \item \textbf{Χείριστη περίπτωση} ($C_{\chi\pi}$) --- το «πιο δύσκολο»
  στιγμιότυπο.
  \item \textbf{Μέση περίπτωση} ($C_{\mu\pi}$) --- σταθμισμένος μέσος πάνω σε όλα
  τα στιγμιότυπα.
\end{itemize}

Πιο τυπικά, αν $D_n$ είναι το σύνολο όλων των στιγμιοτύπων διάστασης $n$:

\[
C_{\beta\pi}(n) = \min_{d \in D_n} \text{κόστος}(d) \qquad \text{(βέλτιστη περίπτωση)}
\]
\[
C_{\chi\pi}(n) = \max_{d \in D_n} \text{κόστος}(d) \qquad \text{(χείριστη περίπτωση)}
\]
\[
C_{\mu\pi}(n) = \sum_{d \in D_n} p(d) \cdot \text{κόστος}(d) \qquad \text{(μέση περίπτωση)}
\]

όπου $p(d)$ η πιθανότητα να εμφανιστεί το στιγμιότυπο $d$ ως είσοδος.

\begin{warningbox}
\textbf{Σχεδόν πάντα στις εξετάσεις χρησιμοποιούμε χείριστη περίπτωση.} Είναι η
εγγύηση --- «ο αλγόριθμος δεν θα πάρει ποτέ παραπάνω από αυτό». Η μέση περίπτωση
χρειάζεται υποθέσεις για την κατανομή της εισόδου, και αυτό σπάνια είναι ξεκάθαρο.
\end{warningbox}

\section{Επεξεργασμένο παράδειγμα: γραμμική αναζήτηση}

\begin{examplebox}
\textbf{Πρόβλημα.} Έχουμε πίνακα $S = (a_1, a_2, \dots, a_n)$ και ζητάμε ένα
στοιχείο $x$. Επιστρέφουμε τη θέση του $x$ στον πίνακα, αν υπάρχει.

\textbf{Αλγόριθμος.}
\begin{codebox}
\begin{verbatim}
i <- 1
Όσο (i != n+1) και (a_i != x):
  i <- i + 1
Εάν i > n: εκτύπωσε("όχι")
Αλλιώς: εκτύπωσε("στοιχείο x στη θέση i")
\end{verbatim}
\end{codebox}

\textbf{Ανάλυση} (βασική στοιχειώδης πράξη: η σύγκριση \texttt{a\_i != x}):

\begin{itemize}
  \item \textit{Βέλτιστη περίπτωση:} το $x$ βρίσκεται στη \textbf{θέση 1}. Μία
  σύγκριση. $C_{\beta\pi}(n) = 1$.
  \item \textit{Χείριστη περίπτωση:} το $x$ είτε δεν υπάρχει, είτε βρίσκεται στη
  \textbf{θέση $n$}. Πρέπει να γίνουν $n$ συγκρίσεις. $C_{\chi\pi}(n) = n$.
  \item \textit{Μέση περίπτωση:} αν υποθέσουμε ομοιόμορφη κατανομή του $x$ στις
  θέσεις $1$ μέχρι $n$, ο αναμενόμενος αριθμός συγκρίσεων είναι
  $\frac{1 + 2 + \dots + n}{n} = \frac{n+1}{2}$, δηλαδή και πάλι γραμμικός στο $n$
  αλλά με σταθερά $1/2$.
\end{itemize}

\textbf{Συμπέρασμα.} Η γραμμική αναζήτηση έχει χείριστη πολυπλοκότητα
\textbf{γραμμική} στο $n$. Η δυαδική αναζήτηση που είδαμε νωρίτερα έχει
\textbf{λογαριθμική} χείριστη πολυπλοκότητα --- γι' αυτό σε μεγάλα δεδομένα η
δεύτερη είναι ασύγκριτα ταχύτερη. Στο επόμενο μάθημα θα μάθουμε να εκφράζουμε αυτή
τη διαφορά τυπικά: $O(n)$ vs $O(\log n)$.
\end{examplebox}

\section{Δεύτερο παράδειγμα: αριθμητικό γινόμενο}

\begin{examplebox}
\textbf{Πρόβλημα.} Δίνονται δύο διανύσματα $a = (a_1, \dots, a_n)$ και
$b = (b_1, \dots, b_n)$. Υπολόγισε το \textbf{αριθμητικό γινόμενο}
$sp = \sum_{k=1}^n a_k \cdot b_k$.

\textbf{Αλγόριθμος.}
\begin{codebox}
\begin{verbatim}
sp <- 0
Για k από 1 έως n:
  sp <- sp + a_k * b_k
Αριθμητικό γινόμενο = sp
\end{verbatim}
\end{codebox}

\textbf{Ανάλυση.} Ο βρόχος εκτελείται ακριβώς $n$ φορές, ανεξάρτητα από την
είσοδο. Κάθε επανάληψη κάνει σταθερό αριθμό πράξεων (1 πολλαπλασιασμό $+$ 1
πρόσθεση $+$ 1 ενημέρωση δείκτη). Άρα:

\[
C_{\beta\pi}(n) = C_{\chi\pi}(n) = C_{\mu\pi}(n) = \Theta(n)
\]

\textbf{Παρατήρηση.} Όταν δεν υπάρχει «εύκολη» ή «δύσκολη» είσοδος --- όταν ο
αλγόριθμος κάνει την ίδια ακριβώς δουλειά για όλα τα στιγμιότυπα της ίδιας
διάστασης --- οι τρεις τύποι πολυπλοκότητας \textbf{συμπίπτουν}.
\end{examplebox}

\section{Το πλάνο του εξαμήνου}

Τέσσερις τεχνικές σχεδιασμού αλγορίθμων, μία θεματική ενότητα γραφημάτων, και μία
ενότητα δομών δεδομένων --- όλα κάθονται πάνω στα θεμέλια που χτίζουμε στα δύο
πρώτα μαθήματα:

\begin{itemize}
  \item \textbf{L01--L02:} Εισαγωγικά \& Ασυμπτωτική ανάλυση (τα θεμέλια).
  \item \textbf{L03--L05:} Divide \& Conquer.
  \item \textbf{L06--L09:} Γραφήματα.
  \item \textbf{L10:} Δομές Δεδομένων.
  \item \textbf{L11--L13:} Άπληστοι αλγόριθμοι (greedy).
  \item \textbf{L14--L17:} Δυναμικός προγραμματισμός.
\end{itemize}

Η ασυμπτωτική ανάλυση είναι \textbf{παντού}: σε κάθε ανάλυση χρόνου που θα
κάνουμε, σε κάθε σύγκριση αλγορίθμων, σε κάθε απόδειξη βελτιστότητας. Γι' αυτό την
παίρνουμε πρώτη.

\section{Βαθμολογία και πρακτικά}

\begin{itemize}
  \item Δύο βαθμοί: \textbf{Πρόοδος (Π)} και \textbf{Γραπτό (Γ)}.
  \item \textbf{Τελικός βαθμός} $= \max\{\Gamma,\ 0.3 \cdot \Pi + 0.7 \cdot
  \Gamma\}$ --- δηλαδή η πρόοδος \textit{βοηθάει}, δεν \textit{βλάπτει}: αν τα πας
  καλά στην πρόοδο και χάλια στο γραπτό, σε «τραβάει» πάνω· αν τα πας καλύτερα στο
  γραπτό, αγνοείται η πρόοδος.
  \item Όριο $40\%$ στο γραπτό για να μετρήσει η πρόοδος.
\end{itemize}

\begin{notebox}
\textbf{Συνιστώμενα προαπαιτούμενα:} Διακριτά Μαθηματικά, Δομές Δεδομένων \&
Τεχνικές Προγραμματισμού. Αν δεν τα έχεις δει --- οι έννοιες που θα χρειαστούμε
χτίζονται μέσα στο μάθημα, αλλά θα πρέπει να επενδύσεις λίγο περισσότερο σε
άσκηση.
\end{notebox}

\section{Βιβλιογραφία}

\begin{itemize}
  \item \textbf{Cormen, Leiserson, Rivest, Stein} --- \textit{Introduction to
  Algorithms} (CLRS). Η «βίβλος» των αλγορίθμων.
  \item \textbf{Kleinberg \& Tardos} --- \textit{Algorithm Design}. Πιο φιλική
  στην απόδειξη ορθότητας, εξαιρετική στο D\&C και το greedy.
  \item \textbf{Dasgupta, Papadimitriou, Vazirani} --- \textit{Algorithms}.
  Συνοπτικότερο, διαθέσιμο ελεύθερα σε pdf.
  \item Συμπληρωματικά: σημειώσεις και διαφάνειες \textbf{Β.~Ζησιμόπουλου}.
\end{itemize}

\begin{recapbox}
\textbf{Τι κρατάμε από αυτή τη διάλεξη}
\begin{enumerate}
  \item Ένας \textbf{αλγόριθμος} είναι μια καλά ορισμένη ακολουθία βημάτων που
  μετασχηματίζει είσοδο σε έξοδο.
  \item Ένα \textbf{πρόβλημα} ορίζει είσοδο/έξοδο. Ένα \textbf{στιγμιότυπο} είναι
  μια συγκεκριμένη είσοδος. Η \textbf{διάσταση} είναι η αριθμητική παράμετρος που
  μετράει το μέγεθος του στιγμιοτύπου --- συνήθως το $n$.
  \item Η \textbf{πολυπλοκότητα} είναι ο αριθμός στοιχειωδών πράξεων ως συνάρτηση
  της διάστασης.
  \item \textbf{Τρεις τύποι}: βέλτιστη, χείριστη, μέση περίπτωση. Σχεδόν πάντα
  χρησιμοποιούμε \textbf{χείριστη}.
  \item Στο μάθημα: 4 τεχνικές σχεδιασμού (D\&C, Greedy, DP) $+$ Γραφήματα $+$
  Δομές Δεδομένων, με την ασυμπτωτική ανάλυση να ενώνει τα πάντα.
\end{enumerate}
\end{recapbox}

\lecturefooter
\end{document}
